\title{Parameter-Efficient Orthogonal Finetuning via Factorization}

\begin{abstract}
Large foundation models have become prevalent, yet the cost of training them from scratch remains prohibitively high. Consequently, finding methods to adapt these sophisticated models efficiently to specific downstream tasks is of paramount importance. This work investigates a rigorous finetuning approach—Orthogonal Finetuning (OFT)—for the purpose of downstream adaptation. While OFT achieves strong generalizability, its reliance on high-dimensional orthogonal matrices results in a substantial number of trainable parameters. To overcome this limitation, we first reinterpret OFT from an information transmission viewpoint and distill several key requirements conducive to parameter efficiency. Taking inspiration from the Cooley-Tukey fast Fourier transform algorithm’s efficient transmission properties, we introduce a compact orthogonal parameterization leveraging butterfly structures. Integrating this factorization into OFT, we put forth a new parameter-efficient finetuning strategy, termed Orthogonal Butterfly~(BOFT). Notably, BOFT generalizes OFT as a special instance within its framework, offering a unified approach to orthogonal finetuning. We further provide comprehensive empirical evaluations, adapting large vision transformers, language models, and text-to-image diffusion architectures to a wide range of downstream tasks in both vision and language domains.
\end{abstract}

\section{Introduction}

Recent developments, exemplified by systems like ChatGPT~\cite{brown2020language,chen2021evaluating} and Stable Diffusion~\cite{rombach2022high}, highlight the extraordinary generalization capabilities of large foundation models. However, the impressive gains in performance have come at the expense of vastly increasing model sizes—GPT-3, for instance, contains approximately 175B parameters. This escalation makes training foundation models from the ground up an increasingly arduous endeavor. As a result, advancing the field largely depends on the capacity to adapt pre-trained foundation models without requiring complete retraining. That is, it is crucial to devise techniques that can efficiently specialize these existing models for downstream applications. Approaches to this challenge generally fall into three categories: \emph{model finetuning}~\cite{hulora2022,zaken2022bitfit,guo2020parameter,raffel2020exploring,qiu2023controlling,zhang2022adaptive,chavan2023one}, which modifies only part of the model’s original parameters while keeping the architecture intact; \emph{adapter tuning}~\cite{rebuffi2017learning,houlsby2019parameter,pfeiffer2020adapterfusion,lin2020exploring,he2021towards}, which integrates additional learnable parameters into the model; and \emph{prompt tuning}~\cite{li2021prefix,lester2021power}, which appends trainable prefix tokens to the input sequence. Among these, model finetuning stands out due to its simplicity, efficacy, and, notably, its lack of added inference latency.

The central idea guiding model finetuning is to maintain a degree of similarity between the original pretrained model and its finetuned counterpart according to certain metrics, thereby preserving the knowledge acquired during pretraining. For example, most contemporary finetuning strategies employ a reduced learning rate; this practical solution helps to limit divergence between the initial and updated models. Considering the inherently structured composition of weight matrices, an alternative, more theoretically motivated method focuses on retaining the relationships within these matrices, specifically by maintaining the pairwise angles between neurons. This observation motivates the development of Orthogonal Finetuning~(OFT)~\cite{qiu2023controlling}, a recent finetuning framework. OFT enforces that neurons within the same layer undergo transformations via a shared orthogonal matrix, which mathematically guarantees the preservation of pairwise neuron angles during finetuning. While OFT has achieved notable improvements in generalization and training stability for text-to-image diffusion models~\cite{qiu2023controlling}, it introduces a significant number of trainable parameters due to the size of the orthogonal matrices involved. To mitigate this challenge, OFT applies a block-diagonal structure, effectively lowering parameter count. Nonetheless, this gain in parameter efficiency comes with notable tradeoffs: the orthogonal matrices are constrained to a fixed sparse block-diagonal form, and the orthogonal operations are limited to acting independently within each block. Although such sparsity saves on resources, it can also create potentially problematic inductive biases (\emph{e.g.}, restricting the range of linear transformations that can be represented and decreasing model flexibility). Furthermore, determining the optimal arrangement for this sparsity structure remains an unresolved problem.

A central challenge lies in constructing a dense orthogonal matrix in a manner that conserves parameters. At first sight, this seems unattainable because representing a dense orthogonal matrix of dimension $d$ conventionally entails $\mathcal{O}(d^2)$ parameters. To circumvent this constraint, our approach innovatively composes a dense orthogonal matrix using a sequence of factorized sparse matrices. The rationale is that it is feasible to economize on the number of trainable parameters by increasing computational complexity, effectively exchanging additional computation for memory savings. Specifically, our method represents the orthogonal matrix as a product of sparse matrices, necessitating repeated multiplications at every iteration during training. To contextualize this matrix factorization, we recast the construction of a dense orthogonal matrix as a problem of information propagation. Within this paradigm, assembling a dense orthogonal matrix through the multiplication of sparse factors can be interpreted as the process of transmitting information over a graph with a grid topology. This transmission-oriented framework provides us with a variety of ways to structure the sparse matrix factorization, thereby managing the parameter count while preserving the capacity to form dense representations.

To optimize parameter efficiency within this information transmission framework, we take inspiration from the butterfly architecture found in the Cooley-Tukey fast Fourier transform algorithm~\cite{cooley1965algorithm}, where nodes communicate in an expedient manner~\cite{klasing1994broadcasting}. Notably, the butterfly graph structure intrinsic to the Cooley-Tukey algorithm yields a systematic method for sparse matrix decomposition, referred to as butterfly factorization. Leveraging this factorization technique, we can express a $d$-dimensional dense matrix as a product of $\mathcal{O}(\log d)$ sparse matrices, each containing $\mathcal{O}(d)$ non-zero elements. By ensuring orthogonality for each sparse factor, our method achieves a highly parameter-efficient orthogonal parameterization, requiring only $\mathcal{O}(d\log d)$ parameters—a marked improvement over the conventional quadratic requirement. Based on this insight, we introduce a new parameter-efficient finetuning strategy: Orthogonal Butterfly~(BOFT). This technique generalizes OFT, serving as an encompassing orthogonal finetuning framework. Both the block-diagonal and butterfly structures have a key commonality—sparsity—which they exploit within orthogonal matrices to decrease the parameter budget. Notably, our strategy offers an instructive comparison with the low-rank structures employed in LoRA~\cite{hulora2022}; while LoRA is based on low-rank matrices, and our work emphasizes sparsity, both approaches belong to prominent classes of structured matrices~\cite{candes2011robust} that deliver parameter efficiency.

In contrast to the block-diagonal configuration adopted by OFT—which negotiates a balance between model expressiveness and structural regularity—BOFT leverages the butterfly structure to facilitate a more seamless transition between matrices drawn from the full orthogonal group (to enhance expressivity) and identity matrices (to promote regularity). This approach allows for the identification of a more compact subset within the orthogonal group that is suitable for specific downstream applications. Given that the butterfly structure underpins many efficient linear transformations, such as the discrete Fourier and discrete cosine transforms, we posit that this structural choice fosters a beneficial inductive bias. This, in turn, is likely to improve model generalization and guard against overfitting. Our reasoning draws on the notable capacity of the butterfly structure to reconstruct various classical linear transforms—something unattainable for the block-diagonal approach in OFT. Our main contributions are outlined below:

\begin{itemize}[leftmargin=*,nosep]
    \item We introduce a new information transmission framework to analyze parameter efficiency in orthogonal finetuning. This perspective recasts the construction of parameter-efficient dense orthogonal matrices as an information transmission task on a grid-structured graph.
    \item Motivated by the butterfly constructions utilized in the Cooley-Tukey algorithm, we present Orthogonal Butterfly, a method for parameter-efficient orthogonal finetuning, integrated within the proposed framework.
    \item We offer theoretical analyses that elucidate why BOFT achieves notable reductions in the number of tunable parameters without sacrificing expressive power or training robustness. Furthermore, due to its matrix factorization, BOFT naturally enables weight interpolation (see Figure~\ref{fig:boft_interpolation}).
    \item For the first time, we demonstrate the applicability of orthogonal finetuning to a wide range of problems outside of controllable text-to-image generation~\cite{qiu2023controlling}. We showcase BOFT across several adaptation tasks in both computer vision and natural language processing, where it significantly outperforms existing leading methods—affirming its efficiency and strong generalization performance.
\end{itemize}

\section{Related Work}

\textbf{Parameter-efficient finetuning (PEFT)}. As foundation models (\emph{e.g.}, \cite{brown2020language,radford2021learning,rombach2022high,kirillov2023segment}) grow in scale and capability, there has been significant momentum towards adapting these models for specific tasks while minimizing the number of tunable parameters~\cite{treviso2023efficient,ding2023parameter,lialin2023scaling}. Among the diverse PEFT strategies~\cite{houlsby2019parameter,aghajanyan2020intrinsic,hulora2022,edalati2022krona,wang2022adamix,gheini2021cross,zaken2022bitfit,guo2020parameter,sung2021training,ansell2022composable,lester2021power,li2021prefix,vu2022spot,he2021towards,mao2021unipelt,karimi2021compacter,liu2022few,sung2022lst,chen2022parameter,jia2022visual,chen2022adaptformer,zhang2022neural,jie2023fact,lian2022scaling,luo2023towards}, reparameterization-based approaches~\cite{aghajanyan2020intrinsic,hulora2022,edalati2022krona} are especially pertinent to our study. LoRA~\cite{hulora2022}, a notable example, refines the pretrained weights by injecting the product of two low-rank matrices, yielding competitive outcomes on various language processing benchmarks. However, LoRA employs a fixed rank across all layers for its added matrices, motivating subsequent innovations~\cite{zhang2022adaptive,valipour2022dylora,zhang2023increlora} that permit adaptive rank selection by layer to optimize parameter allocation. The straightforwardness of this low-rank reparameterization has led to its widespread adoption~\cite{dettmers2023qlora,zi2023delta,chavan2023one}. Drawing from research on hyperspherical energy and generalization~\cite{liu2018learning,liu2021orthogonal}, \cite{qiu2023controlling} introduces orthogonal finetuning (OFT), an alternative strategy for tuning text-to-image diffusion models. OFT modifies neuron activations in a layer using an orthogonal transformation, leading to more robust generalization and training stability when compared to LoRA. Nevertheless, OFT generally incurs a greater parameter count than LoRA, motivating efforts to enhance its parameter efficiency. Additionally, it remains unclear if OFT can be extended to other adaptation challenges outside of text-to-image diffusion models~\cite{qiu2023controlling}. BOFT addresses the parameter efficiency limitation in OFT by employing butterfly factorization, thereby enabling a broader evaluation of orthogonal finetuning in general adaptation scenarios.

\textbf{Butterfly structures}. The radix-2 Cooley-Tukey method achieves an $N$-point discrete Fourier transform by recursively decomposing the problem into smaller $\frac{N}{2}$-point transforms. This decomposition gives rise to a characteristic butterfly architecture, represented as a product of multiple sparse matrices (often referred to as a butterfly matrix). Butterfly matrices have previously served in parametrizing orthogonal matrices to eliminate the need for pivoting in Gaussian elimination and to enhance computational efficiency~\cite{parker1995random}, in facilitating stable training in recurrent neural networks~\cite{jing2017tunable}, and for efficient kernel approximations~\cite{munkhoeva2018quadrature}. Other works~\cite{dao2019learning,dao2019kaleidoscope} have explored butterfly parameterizations for fast linear transformation learning. The efficient training of neural networks through butterfly matrices has been pursued in~\cite{chen2022pixelated,dao2022monarch}. Additionally, butterfly structures have been leveraged for accelerated matrix-vector computations~\cite{michielssen1996multilevel,o2010algorithm}, compact matrix representation~\cite{li2015butterfly}, and improvements in network data transmission~\cite{klasing1994broadcasting,li2003linear,rajasingh2016transmission}. Departing from these prior applications, our work leverages butterfly structures specifically to improve the parameter efficiency of OFT.

\section{An Information Transmission View on Orthogonal Finetuning}

We begin by outlining essential background on OFT. In order to adapt the pretrained weight matrix, OFT introduces a reparameterization where the updated weights are expressed as a product of a learnable orthogonal matrix and the fixed, pretrained weight matrix. Unlike LoRA, which incorporates parameter updates through the addition of a low-rank matrix, OFT employs a multiplicative orthogonal matrix for its weight modifications. LoRA's parameter-efficiency is achieved via a low-rank configuration, whereas OFT's initial design leverages a block-diagonal orthogonal structure~\cite{qiu2023controlling}—notably, reducing the diagonal block size increases the parameter-efficiency of OFT. A direct visual comparison is shown in Figure~\ref{fig:comp_lora}. The key rationale for using orthogonal transformations in finetuning is to retain the angular relationships between neurons~\cite{liu2018learning,liu2021orthogonal,qiu2023controlling}, thereby largely maintaining the pretrained model’s semantic information.

Formally, OFT learns an orthogonal matrix $\bm{R}\in\mathbb{R}^{d\times d}$ associated with a given pretrained linear weight $\bm{W}^0\in\mathbb{R}^{d\times n}$, and alters the forward computation from $\bm{z}=(\bm{W}^0)^\top\bm{x}$ to $\bm{z}=(\bm{R}\bm{W}^0)^\top\bm{x}$, where $\bm{x}\in\mathbb{R}^{d}$ is the input and $\bm{z}\in\mathbb{R}^{n}$ is the output vector. For zero initialization, $\bm{R}$ is set as the identity matrix. To maintain the orthogonal property of $\bm{R}$ throughout finetuning, we adopt the Cayley parameterization as described in~\cite{qiu2023controlling,liu2021orthogonal}: $\bm{R}=(\bm{I}+\bm{Q})(\bm{I}-\bm{Q})^{-1}$, where $\bm{Q}$ is skew-symmetric, i.e., $\bm{Q}=-\bm{Q}^\top$. In order to improve parameter-efficiency, a block-diagonal form is introduced by expressing $\bm{R}$ as $\text{diag}(\bm{R}_1,\bm{R}_2,\cdots,\bm{R}_r)$, where each $\bm{R}_i\in\mathbb{R}^{b\times b}$ for all $i$ is an orthogonal matrix and $br=d$. The benefit of parameter-efficiency with the block-diagonal structure is counterbalanced by the imposed constraint: this approach implicitly assumes that the dimensions of each neuron (i.e., columns of $\bm{W}^0$) are split into $r$ groups, with each group undergoing its own distinct orthogonal transformation. Even though the block-diagonal design is empirically robust, splitting neuron dimensions solely by index lacks conceptual justification, which points to the advantages of using dense orthogonal matrices. This consideration leads to a pertinent research question: \emph{Is it possible to create a dense orthogonal matrix while preserving parameter-efficiency?}

To investigate this problem, we suggest constructing a dense orthogonal matrix as a product of several sparse orthogonal matrices. In order to analyze the sparsity structure of orthogonal matrix factorization from a consistent yet accessible viewpoint, we recast the task of forming a dense orthogonal matrix in OFT as analogous to an information transmission scenario. More concretely, synthesizing a dense matrix $\bm{R}\in\mathbb{R}^{d\times d}$ by sequentially multiplying $m$ square matrices $\thickmuskip=2mu \medmuskip=2mu \bm{R}=\bm{B}_m\bm{B}_{m-1}\cdots\bm{B}_1$ can be interpreted as transmitting signals through a grid comprising $d\times (m+1)$ nodes, as depicted in Figure~\ref{fig:info_trans}. The impetus for this information transmission perspective arises from viewing a $d$-dimensional dense square matrix as representing complete interconnections from one set of $d$ nodes to another. For instance, in the matrix $\bm{R}$, a zero at position $R_{ij}$ signifies the impossibility of passing information from node $j$ to node $i$, whereas a nonzero entry denotes the ability for information to be transferred. Consequently, the decomposition of $\bm{R}$ into several factors $\bm{B}_m\bm{B}_{m-1}\cdots\bm{B}_1$ may be seen as a process of successive information transmission dictated by the network structure induced by each $\bm{B}_i$. The information first traverses the graph associated with $\bm{B}_1$, proceeding sequentially up to $\bm{B}_m$. As a tangible illustration, Figure~\ref{fig:info_trans} considers the decomposition $\bm{R}=\bm{B}_5\bm{B}_4\bm{B}_3\bm{B}_2\bm{B}_1$, along with visualizations of their respective sparsity patterns and induced graph. The depicted graph results from unwinding the sequence of matrix multiplications. Within this induced network, each matrix $\bm{B}_i$ serves as a connectivity blueprint linking nodes from the $i$-th layer to those in the $(i+1)$-th layer; specifically, the $(j_1,j_2)$ entry of $\bm{B}_i$ encodes whether a directed edge exists from the $j_2$-th node at level $i$ to the $j_1$-th node at level $(i+1)$ (a zero meaning no edge). For the product $\bm{B}_5\bm{B}_4\bm{B}_3\bm{B}_2\bm{B}_1$ to yield a dense result, it is necessary that every node at the initial level is capable of conveying information to each node at level six. If we consider instead the partial product $\bm{R}=\bm{B}_4\bm{B}_3\bm{B}_2\bm{B}_1$, corresponding to information flow from source nodes at level one to destination nodes at level five, we see, for example, that node $1$ is unable to reach node $3$, which is mirrored by a zero at location $(3,1)$ in the associated sparsity pattern. This correspondence between reachability in the induced network and the sparsity pattern persists for $\bm{R}=\bm{B}_3\bm{B}_2\bm{B}_1$ and $\bm{R}=\bm{B}_2\bm{B}_1$. More generally, achieving a dense $\bm{R}\in\mathbb{R}^{d\times d}$ requires the $m$ factor matrices $\bm{B}_m,\cdots,\bm{B}_1$ to collectively encode a set of directed edges within a $d\times (m+1)$ grid such that each edge links two adjacent layers (i.e., between neighboring columns), and any node in the initial layer can reach all nodes in the final, $(m+1)$-th, layer through some sequence of edges.

Viewing the problem through the lens of information transmission yields valuable insights into the sparsity pattern present in the factorization matrices used in OFT. In Figure~\ref{fig:blk_diag}, the block-diagonal arrangement of $\bm{R}$ in the standard OFT is depicted. This configuration reduces the trainable parameter count; however, it prevents $\bm{R}$ from being a dense matrix. Our objective, therefore, is to construct a dense orthogonal matrix by composing $m$ sparse orthogonal matrices, aiming to minimize the number of required trainable parameters. From the information transmission perspective, two key design criteria emerge: \emph{(i)} \textbf{dense connectivity}, where each node in the initial layer maintains at least one connection to all nodes in the final layer, and \emph{(ii)} \textbf{minimum free edges}, requiring that the sum of all edges be as low as possible without violating the orthogonality constraint. 

Orthogonality, in this context, imposes specific limitations on connections between consecutive levels. For instance, each $\bm{B}_i$ must be of full rank—a prerequisite for orthogonality—which necessitates at least $d$ edges establishing a one-to-one correspondence between nodes in the $i$-th and $(i+1)$-th levels. Thus, the minimum number of edges between levels is $d$ (as shown, for example, in Figure~\ref{fig:info_trans} with 4 edges). These edges are dictated by the requirement for orthogonality and are not tunable parameters (e.g., in a $d\times d$ orthogonal matrix, the $d$ nonzero entries must be exactly $\pm 1$). The orthogonality condition therefore leads to greater parameter efficiency compared to unrestricted dense matrices, but also introduces dependencies between different edges: for example, in a $2\times 2$ orthogonal matrix, if the $(1,1)$ entry is zero, the $(2,2)$ entry must be zero as well, so omitting one edge enforces the absence of another. Consequently, depending on the edge configuration, orthogonality might require the inclusion or removal of certain edges.

Analyzing the pattern of nonzero entries in a composed orthogonal matrix highlights the value of the information transmission framework in OFT: it enables focus exclusively on the locations of nonzero, trainable elements within $\bm{R}$, independent of their actual magnitudes.

A straightforward fully connected approach between two layers entails $\mathcal{O}(d^2)$ connections (essentially, one full dense orthogonal matrix), resulting in $d^2 - d$ connections (for $d = 4$, this amounts to 12 connections). Figure~\ref{fig:info_trans} illustrates a possible matrix decomposition that utilizes a total of $10$ connections, which is notably fewer than required by a single dense orthogonal matrix. This perspective provides a graphical framework for analyzing the parameter-efficiency of OFT, facilitating the development of viable matrix decompositions. Motivated by a distinctive topology employed in the Cooley-Tukey algorithm—namely, butterfly graphs~\cite{cooley1965algorithm}—we observe that these structures are capable of densely connecting $d$ inputs and $d$ outputs with only $\mathcal{O}(d\log d)$ edges. For example, the topology displayed in Figure~\ref{fig:info_trans} necessitates $10$ edges to achieve complete connectivity, whereas the butterfly architecture accomplishes this with just $8$ edges. In what follows, we detail how leveraging the butterfly structure leads to more efficient parameterizations.

\section{Orthogonal Parameterization by Butterfly Factorization}

Initially developed for the Cooley-Tukey algorithm to facilitate the fast Fourier transform, the butterfly structure possesses the ability to transmit localized modifications in the frequency domain into global transformations in the spatial domain. This quality parallels the information flow challenge we address: ensuring that each node at the initial layer is able to communicate with all nodes at the terminal layer. Moreover, butterfly networks have been widely adopted as efficient topologies in computer networking~\cite{solihin2015fundamentals,li2011linear} to enable rapid data transmission. Considering $k \geq 2$ as a power of $2$, we introduce the \emph{butterfly factor} $\bm{B}^F(k)$ as follows:

\begin{equation}\label{eq:but_factor}
\begin{aligned}
    \bm{B}^F(k) = \begin{bmatrix}
    \text{diag}(\bm{d}_1) & \text{diag}(\bm{d}_2) \\
    \text{diag}(\bm{d}_3) & \text{diag}(\bm{d}_4)
    \end{bmatrix} \in \mathbb{R}^{k \times k},
\end{aligned}
\end{equation}

where each $\bm{d}_i \in \mathbb{R}^{\frac{k}{2}}$ for $i=1,2,3,4$. Assuming $d = 2^N$, we recursively define the $d$-dimensional \emph{butterfly matrix} $\bm{B}(d) \in \mathbb{R}^{d \times d}$ as

\begin{equation}
\begin{aligned}
    \bm{B}(d) = \tilde{\bm{B}}(d, d) \cdot \begin{bmatrix}
    \bm{B}_1(\frac{d}{2}) & \bm{0} \\
    \bm{0} & \bm{B}_2(\frac{d}{2})
    \end{bmatrix} = \tilde{\bm{B}}(d, d) \tilde{\bm{B}}(d, \frac{d}{2}) \cdots \tilde{\bm{B}}(d, 2),
\end{aligned}
\end{equation}

Here, $\bm{B}_1(\frac{d}{2})$ and $\bm{B}_2(\frac{d}{2})$ represent two butterfly matrices of dimension $\frac{d}{2}$. Next, we introduce the notion of the \emph{butterfly component}, formulated as $\thickmuskip=2mu \medmuskip=2mu \tilde{\bm{B}}(d,k)=\text{diag}(\bm{B}^F_1(k),\cdots,\bm{B}^F_{d/k}(k))$. This object is a block-diagonal matrix with dimension $d\times d$ in which each block has size $k$, where each $\bm{B}^F_i(k)$ (for all $i$) corresponds to a butterfly factor as specified in Equation~\ref{eq:but_factor}. The butterfly matrix provides a way to represent an orthogonal matrix by parameterization. This is achieved by ensuring that every multiplicative component $\tilde{\bm{B}}(d,k)$ within the overall butterfly matrix $\bm{B}(d)$ retains orthogonality. To illustrate, consider the block-diagonal matrix $\tilde{\bm{B}}(d,2)$, composed of $2 \times 2$ blocks; orthogonality can be readily imposed on each block via the Cayley transform or equivalently, via $2$-dimensional rotations, as in \cite{liu2021orthogonal,qiu2023controlling}. Since the non-zero structure in each butterfly component can be re-expressed as a permutation of the pattern in $\tilde{\bm{B}}(d,2)$, all such components are amenable to parameterization as orthogonal matrices. This approach leads to an efficient orthogonal matrix parameterization based on $2\times2$ orthogonal submatrices. Following the generalization in \cite{chen2022pixelated}, we define the block butterfly component $\tilde{\bm{B}}^b(d,k)$, where each element in $\bm{d}_i$ becomes a $b \times b$ matrix. Orthogonality of $\tilde{\bm{B}}^b(d,2)$ is maintained by parameterizing each $2b\times2b$ block as orthogonal. For larger $k$ $(k>2)$, the non-zero pattern in the butterfly components $\tilde{\bm{B}}^b(d,k)$ corresponds to a block permutation of that in $\tilde{\bm{B}}^b(d,2)$, so these components can also be orthogonally parameterized in a similar fashion. Putting these elements together, the forward propagation in BOFT is expressed as follows:

\begin{equation}
    \begin{aligned}\nonumber
    \bm{z}=\big(\bm{R}(m,b)\cdot\bm{W}^0\big)^\top\bm{x},~~~~\text{s.t.}~\bigg{\{}\bm{R}(m,b)=\prod_{i=1}^m \tilde{\bm{B}}^b_{(d,i)}~~\&~~\underbrace{\big(\tilde{\bm{B}}^b_{(d,j)}\big)^\top\tilde{\bm{B}}^b_{(d,j)}=\tilde{\bm{B}}^b_{(d,j)}\big(\tilde{\bm{B}}^b_{(d,j)}\big)^\top\!=\bm{I}_d}_{\forall j\in [1, m]}\bigg{\}},
    \end{aligned}
\end{equation}

Here, for clarity, we use the shorthand $\tilde{\bm{B}}^b_{(d,i)}$ to represent $\tilde{\bm{B}}^b(d,2^{m-i+1})$, and $\bm{I}_d$ is the $d \times d$ identity matrix. The matrix $\bm{R}(m,b)\in\mathbb{R}^{d\times d}$ is constructed as the product of several orthogonal butterfly matrices. For succinct notation, we write BOFT using $\bm{R}(m,\frac{b}{2})$ as BOFT($m$,$b$), where $b\geq2$. When $\thickmuskip=2mu \medmuskip=2mu m=1$, the formulation BOFT($1$,$b$) recovers the block-diagonal OFT~\cite{qiu2023controlling} with a block size $b$. Choosing BOFT($1$,$d$) gives the standard OFT~\cite{qiu2023controlling} where the orthogonal matrix is unconstrained and fully dense. For BOFT($\log \frac{2d}{b}$,$b$), we employ the block butterfly structure $\bm{B}^{\frac{b}{2}}(d)$ for $\bm{R}$, which results in a dense orthogonal matrix. Typically, BOFT($m$,$b$) introduces $\frac{1}{2}(b-1)dm$ trainable parameters when adapting a linear transformation of dimensions $d \times n$. In the special case where the butterfly structure is used, namely $m=\log d, b=2$, BOFT reduces the parameter count to $\mathcal{O}(d\log d)$. This is more efficient compared to the original OFT, which requires $\mathcal{O}(d^2)$ parameters for a fully dense orthogonal matrix, and the block-diagonal variant with $r$ blocks demands $\mathcal{O}(bd)$. Consequently, whereas original OFT requires setting $b=d$ to form a dense orthogonal structure, BOFT is able to do so for any choice of $b$.

\textbf{Identity Initialization for BOFT}. It is common in finetuning strategies to start from the original pretrained weights to ensure the adapted model remains close to the initial parameters. LoRA, for instance, initializes its low-rank updates with zeros. In the case of BOFT, we initialize each butterfly block as an identity matrix (that is, the associated skew-symmetric matrices in the Cayley transform start as zero matrices).

\textbf{Multiplicative Dropout for BOFT}. In LoRA~\cite{hulora2022}, Dropout is applied to the low-rank modifications to mitigate overfitting. Standard Dropout~\cite{srivastava2014dropout} integrates naturally with LoRA, but is not directly suitable for BOFT due to its multiplicative nature in updating weights. To overcome this, we introduce a multiplicative Dropout scheme specific to BOFT. Since $\bm{R}(m,b)$ consists of $m$ butterfly matrices, which can be rearranged as $2b\times 2b$ block-diagonal orthogonal matrices, multiplicative Dropout involves two steps: selecting randomly $p_1$ percent of butterfly matrices, and within each, masking $p_2$ percent of diagonal blocks, replacing the chosen blocks by identity matrices.

\section{Intriguing Insights and Discussions}

\textbf{Expressivity of BOFT}. While the butterfly architecture combined with permutations has been shown to exactly recover several well-known fast linear transforms~\cite{dao2019learning,dao2019kaleidoscope} (\emph{e.g.}, fast Fourier transform, Hadamard transform), its ability to approximate arbitrary orthogonal matrices using our orthogonal butterfly matrices remains an open question. To investigate this, we perform a simulation in which a randomly generated dense orthogonal matrix~\cite{anderson1987generation} of dimensions $512\times 512$ is approximated. The results, visualized in Figure~\ref{fig:para_eff}, are averaged over 10 independent initializations. The y-axis measures the approximation error, whereas the x-axis reflects the number of trainable parameters used. Each curve (color-coded) corresponds to BOFT models with a particular block size; the leftmost points depict the performance of BOFT($1$,$b$), which refers to the basic block-diagonal OFT with block size $b$. Empirically, BOFT attains greater parameter efficiency than OFT. For instance, BOFT($9$,$2$) is more expressive than BOFT($1$,$16$) despite using significantly fewer parameters. More generally, increasing $m$ and decreasing $b$ tends to increase parameter efficiency: BOFT($6$,$4$) achieves an approximation error comparable to BOFT($2$,$16$) but with fewer parameters. This suggests that the butterfly matrix captures a more restrictive subset of the orthogonal group than block-diagonal forms do, leading to the provable expressiveness advantage of BOFT over OFT for a fixed block size.

\begin{theorem}[Expressivity of BOFT]\label{thm:exp_boft}
    For equal block sizes, BOFT is strictly more expressive than OFT. By forming products of butterfly matrices as $\bm{B}_{d-1,1}(d)\bm{B}^\top_{d-1,2}(d)\cdots\bm{B}_{1,1}(d)\bm{B}^\top_{1,2}(d)$, where each $\bm{B}_{i,j}(d)$ denotes a butterfly matrix, it is possible to approximate any $d\times d$ orthogonal matrix.
\end{theorem}

Theorem~\ref{thm:exp_boft} provides a natural way to generalize BOFT by constructing the final orthogonal matrix as $\thickmuskip=2mu \medmuskip=2mu \bm{R}^G(m_1,b_1,m_2,b_2,l)=\bm{R}_{l,1}(m_1,b_1)\bm{R}_{l,2}^T(m_2,b_2)\cdots\bm{R}_{1,1}(m_1,b_1)\bm{R}_{1,2}^T(m_2,b_2)$, where $\bm{R}_{i,j}^T(m,b)$ are the orthogonal matrices arising in the BOFT formulation. If we set $m_1=m_2=\log d$, $b_1=b_2=2$, and $l=d-1$, the resulting $\bm{R}^G(m_1,b_1,m_2,b_2,l)$ spans the entire orthogonal group—this construction is also known as the kaleidoscope hierarchy~\cite{dao2019kaleidoscope}. However, it is important to emphasize that greater expressivity does not inherently result in superior finetuning performance: for example, full finetuning, despite being universally expressive, often produces suboptimal outcomes. Thus, striking a balance between expressive power and structural regularization is critical for effective model generalization during finetuning. BOFT augments the parameter space accessible in finetuning via structured priors, thereby enabling an improved balance between expressivity and regularity.

\textbf{Spectral properties}. Compared to LoRA, orthogonal finetuning typically demonstrates superior preservation of spectral characteristics, as it exactly maintains the spectral norm of the original pretrained weight matrix $\bm{W}^0$. This can be demonstrated via singular value decomposition: $\bm{W}^0=\bm{U}\bm{\Sigma}\bm{V}^\top$, with $\bm{U},\bm{V}$ being orthogonal matrices and $\bm{\Sigma}$ containing the singular values along its diagonal. Both OFT and BOFT operate by left-multiplying an orthogonal matrix $\bm{R}$, resulting in the adjusted weights $\bm{R}\bm{U}\bm{\Sigma}\bm{V}^\top$. This process preserves the largest singular value, i.e., the spectral norm of $\bm{W}^0$, without alteration. The ability to retain this property has been evidenced to significantly enhance both the stability of model training and generalization capabilities~\cite{miyato2018spectral,yoshida2017spectral}. Additional mathematical insights are detailed in Appendix~\ref{app:sn_property}.

\textbf{Orthogonal finetuning as learning bilinear similarity}. The BOFT approach can be formulated as an optimization problem where the bilinear similarity $\bm{w}^0_i\bm{R}\bm{x}$ is learned, with $\bm{w}^0_i$ denoting the $i$-th neuron (that is, the column of the weight matrix $\bm{W}^0$). In this framework, BOFT interprets the learning process as identifying a bilinear similarity matrix $\bm{R}$ under a stringent orthogonality constraint, thereby linking the approach closely to concepts found in distance metric learning~\cite{xing2002distance} and bilinear forms~\cite{roman2005advanced}.

\textbf{Inductive bias and generalization in BOFT}. Within BOFT, the matrix $\bm{R}(m,b)$ generally belongs to a structured subset of the orthogonal group, which places restrictions on the range of models considered and inherently imposes an inductive bias. We posit that this type of structured inductive bias, especially as introduced via butterfly factorization, enhances generalization. This is due to its resemblance to the organized patterns present in numerous well-established linear transformations~\cite{dao2019kaleidoscope}, such as the discrete Fourier transform, discrete sine/cosine transform, and Hadamard transform. Furthermore, the incorporation of sparse matrix factorization within BOFT may impart additional, implicit inductive biases~\cite{gunasekar2017implicit,li2018algorithmic,liu2019neural}.

\textbf{Contrast with butterfly-based sparse training}. Prior research~\cite{dao2019learning,dao2019kaleidoscope,chen2022pixelated,dao2022monarch} has extensively explored the use of butterfly parameterization for sparse training. These studies primarily emphasize the direct reparameterization of weight matrices using butterfly structures and training neural networks from initialization. Notably, \cite{dao2022monarch} explores finetuning by first projecting pretrained weights onto a form of butterfly matrices, followed by optimizing these projections for specific downstream applications. In contrast, BOFT introduces a novel finetuning approach, utilizing layer-shared weight matrices to transform model weights.

\input{rebuttal-results}

\section{Concluding Remarks and Limitations}

This work introduces Orthogonal Butterfly, a broadly applicable and parameter-efficient method for finetuning foundation models through the use of butterfly architectures. The principal insight driving parameter-efficiency in our approach is to express a dense orthogonal matrix as a product of multiple sparse orthogonal matrices. To systematically identify such matrix decompositions, we devise a graphical information transmission perspective. Using this approach, we find that the butterfly structure aptly serves our goal of sparse orthogonal matrix factorization. Extensive experiments highlight the practical utility of BOFT for finetuning across large-scale language models, vision models, and text-to-image generative systems. The results further support BOFT's capabilities as a versatile finetuning technique.

Nevertheless, BOFT exhibits certain limitations. The composition of multiple orthogonal matrices in BOFT results in moderately increased training time overhead relative to OFT. Enhancing the training efficiency of BOFT thus remains an open research question. On the positive side, once finetuning concludes, the orthogonal matrices learned by BOFT can be consolidated into the pretrained model, incurring no extra inference-time cost. Additionally, it remains uncertain whether the butterfly architecture is the most optimal for information transmission. Our proposed information transmission viewpoint offers opportunities to incorporate ideas from fields such as computer networking, where network topologies for efficient data dissemination are rigorously studied. This opens possibilities for leveraging alternative, potentially more efficient structures for constructing dense orthogonal matrices.